% !TeX root = ../../infdesc.tex

\subsection*{Finite sets}

\begin{chapex}
\label{cqFunctionInjectiveIffSurjective}
Let $n \in \mathbb{N}$ and let $f : [n] \to [n]$ be a function. Prove that $f$ is injective if and only if $f$ is surjective.
\hint{cqFunctionInjectiveIffSurjective}{%
Proceed by induction on $n$. The proof of \Cref{thmJectionsAndSizeOfNaturalNumbers} might provide some useful inspiration.
}
\end{chapex}

\begin{chapex}
\label{cqSizeOfZModN}
Prove that $|\mathbb{Z}/n\mathbb{Z}|=n$ for all integers $n > 0$.
\hint{cqSizeOfZModN}{%
There is a particularly natural function $f : [n] \to \mathbb{Z}/n\mathbb{Z}$ that you should define, and then prove that $f$ is a bijection.
}
\end{chapex}

In \Crefrange{cqFiniteSymmetricDifferenceBegin}{cqFiniteSymmetricDifferenceEnd} the symbol $\triangle$ denotes the\textit{symmetric difference} operation defined in \Cref{defSymmetricDifference}.

\begin{chapex}
\label{cqFiniteSymmetricDifferenceBegin}
\label{cqSizeOfSymmetricDifference}
Let $A$ and $B$ be sets.
\begin{enumerate}[(a)]
\item Prove that if $A$ and $B$ are finite, then $A \triangle B$ is finite, and express $|A \triangle B|$ in terms of $|A|$, $|B|$ and $|A \cap B|$.
\item Prove that if $A \triangle B$ is finite then $A \setminus B$ and $B \setminus A$ are finite.
\item Prove that there are infinite sets $A$ and $B$ such that $A \triangle B$ is finite.
\end{enumerate}
\hint{cqSizeOfSymmetricDifference}{%
The result of \Cref{cqSymmetricDifferenceInTermsOfOtherSetOperations} is helpful ffor (a) and (b), together with the results proved in \Cref{secFiniteSets}. The result of \Cref{cqSymmetricDifferenceOfSetWithItselfAndEmptySet} is helpful in constructing an example for part (c).
}
\end{chapex}

\begin{chapex}
\label{cqFiniteSymmetricDifferenceEnd}
\label{cqSymmetricDifferenceTwoOfThree}
Let $A$, $B$ and $C$ be sets. Prove that if $A \triangle B$ and $B \triangle C$ are finite, then $A \triangle C$ is finite, where the notation $\triangle$ refers to the \textit{symmetric difference} operation defined in \Cref{defSymmetricDifference}.
\hint{cqSymmetricDifferenceTwoOfThree}{%
Start by proving that $(A \triangle C) \cap B$ and $(A \triangle C) \setminus B$ are both finite, using the definitions of the set operations involved and the given hypotheses.
}
\end{chapex}

\subsection*{True--False questions}

\tfquestiontext{cqFiniteInfiniteSetsTFBegin}{cqFiniteInfiniteSetsTFEnd}

\begin{chapex} % False
\label{cqFiniteInfiniteSetsTFBegin}
Every finite set has a unique enumeration.
\end{chapex}

\begin{chapex} % True
The union of a finite number of finite sets is finite.
\end{chapex}

\begin{chapex} % True
Every finite set is a proper subset of another finite set.
\end{chapex}

\begin{chapex} % False
\label{cqFiniteInfiniteSetsTFEnd}
There is a finite set with infinitely many subsets.
\end{chapex}

\subsection*{Always--Sometimes--Never questions}

\asnquestiontext{cqFiniteInfiniteSetsASNBegin}{cqFiniteInfiniteSetsASNEnd}

\begin{chapex} % Sometimes
\label{cqFiniteInfiniteSetsASNBegin}
Let $X$ and $Y$ be finite sets and let $f : X \to Y$ be a function. If $|X| \le |Y|$, then $f$ is injective.
\end{chapex}

\begin{chapex} % Never
Let $X$ be a finite set. Then there is some proper subset $U \subsetneqq X$ such that $|U| = |X|$.
\end{chapex}

\begin{chapex} % Sometimes
Let $X$ be an infinite set. Then $\mathbb{N} \subseteq X$.
\end{chapex}

\begin{chapex} % Always
Let $X$ be a finite set. Then $\mathcal{P}(X)$ is finite.
\end{chapex}

\begin{chapex} % Never
Let $a,b \in \mathbb{R}$ with $a<b$. Then the set $[a,b] \cap \mathbb{Q}$ is finite.
\end{chapex}

\begin{chapex} % Sometimes
Let $X$ and $Y$ be sets with $X$ finite, and let $f : X \to Y$ be an injection. Then $Y$ is finite.
\end{chapex}

\begin{chapex} % Always
\label{cqFiniteInfiniteSetsASNEnd}
Let $X$ and $Y$ be sets with $X$ finite, and let $f : X \to Y$ be a surjection. Then $Y$ is finite.
\end{chapex}